\section{Properties of Diffusion Planners}
\label{sec:properties}

We discuss a number of \method's important properties, focusing on those that are are either distinct from standard dynamics models or unusual for non-autoregressive trajectory prediction.

% \method{} has a number of properties that distinguish it from existing single-step autoregressive dynamics models commonly used in deep model-based reinforcement learning: 

\textbf{Learned long-horizon planning.}~~
Single-step models are typically used as proxies for ground-truth environment dynamics $\vf$, and as such are not tied to any planning algorithm in particular.
In contrast, the planning routine in Algorithm~\ref{alg:rl} is closely tied to the specific affordances of diffusion models.
Because our planning method is nearly identical to sampling (with the only difference being guidance by a perturbation function $h(\btau{})$), \method's effectiveness as a long-horizon predictor directly translates to effective long-horizon planning.
We demonstrate the benefits of learned planning in a goal-reaching setting in \textbf{Figure~\ref{fig:properties}a}, showing that \method is able to generate feasible trajectories in the types of sparse reward settings where shooting-based approaches are known to struggle.
We explore a more quantitative version of this problem setting in Section~\ref{sec:maze2d}.

\vspace{.2em}
\textbf{Temporal compositionality.}~~
Single-step models are often motivated using the Markov property, allowing them to compose in-distribution transitions to generalize to out-of-distribution trajectories.
Because \method generates globally coherent trajectories by iteratively improving local consistency (Section~\ref{sec:model}), it can also stitch together familiar subsequences in novel ways.
In \textbf{Figure~\ref{fig:properties}b}, we train \method on trajectories that only travel in a straight line, and show that it can generalize to v-shaped trajectories by composing trajectories at their point of intersection.

\vspace{.2em}
\textbf{Variable-length plans.}~~
Because our model is fully convolutional in the horizon dimension of its prediction, its planning horizon is not specified by architectural choices.
Instead, it is determined by the size of the input noise $\btau{N} \sim \mathcal{N(\mathbf{0}, \mathbf{I})}$ that initializes the denoising process, allowing for variable-length plans (\textbf{Figure~\ref{fig:properties}c}).

\vspace{.2em}
\textbf{Task compositionality.}~~
While \method contains information about both environment dynamics and behaviors, it is independent of reward function.
Because the model acts as a prior over possible futures, planning can be guided by comparatively lightweight perturbation functions $h(\btau{})$ (or even combinations of multiple perturbations) corresponding to different rewards.
We demonstrate this by planning for a new reward function unseen during training of the diffusion model (\textbf{Figure~\ref{fig:properties}d}).

%% generate with:
%%      python plotting/read_results.py {L9 mode = 'maze2d'}
%%      on branch models-antmaze-temporal

\begin{table}[t]
\centering
\small
\begin{tabular*}{\columnwidth}{@{\extracolsep{\fill}}lrrrrr}
\toprule
\multicolumn{1}{c}{\bf Environment} &
    \multicolumn{1}{c}{\bf MPPI} &
    \multicolumn{1}{c}{\bf CQL} &
    \multicolumn{1}{c}{\bf IQL} &
    \multicolumn{1}{c}{\bf \method} \\
\midrule
Maze2D~~~~U-Maze &
    $33.2$ &
    $5.7$ &
    $47.4$ &
    \highlight{\color{highlight}113.9} \scriptsize{\raisebox{1pt}{$\pm 3.1$}} \\ 
Maze2D~~~~Medium &
    $10.2$ &
    $5.0$ &
    $34.9$ &
    \highlight{\color{highlight}121.5} \scriptsize{\raisebox{1pt}{$\pm 2.7$}} \\ 
Maze2D~~~~Large &
    $5.1$ &
    $12.5$ &
    $58.6$ &
    \highlight{\color{highlight}123.0} \scriptsize{\raisebox{1pt}{$\pm 6.4$}} \\
\midrule
\multicolumn{1}{c}{\bf Single-task Average} & 16.2 & 7.7 & 47.0 & \highlight{\color{highlight}119.5} \hspace{.58cm} \\
\vspace{-.2cm} \\
%%%%%%%%
\toprule
Multi2D~~~~U-Maze & 
    $41.2$ &
    - &
    $24.8$ &
    \highlight{\color{highlight}128.9} \scriptsize{\raisebox{1pt}{$\pm 1.8$}} \\ 
Multi2D~~~~Medium &
    $15.4$ &
    - &
    $12.1$ &
    \highlight{\color{highlight}127.2} \scriptsize{\raisebox{1pt}{$\pm 3.4$}} \\ 
Multi2D~~~~Large &
    $8.0$ &
    - &
    $13.9$ &
    \highlight{\color{highlight}132.1} \scriptsize{\raisebox{1pt}{$\pm 5.8$}} \\
\midrule
\multicolumn{1}{c}{\bf Multi-task Average} & 21.5 & - & 16.9 & \highlight{\color{highlight}129.4} \hspace{.58cm} \\ 
\bottomrule
\end{tabular*}
% \vspace{-5pt}
\caption{
    \textbf{(Long-horizon planning)}
    The performance of \method{} and prior model-free algorithms in the Maze2D environment, which tests long-horizon planning due to its sparse reward structure.
    The Multi2D setting refers to a multi-task variant with goal locations resampled at the beginning of every episode.
    \method{} substantially outperforms prior approaches in both settings.
    Appendix~\ref{app:baselines} details the sources for the scores of the baseline algorithms.
}
% \vspace{-.1cm}
\label{table:maze2d}
% \vspace{-15pt}
\end{table}

